\section{Source analysis and the new contributions}

\subsection{The pinned source and its mathematical scope}

The source is Vladimir Reshetnikov's ProveIt repository. We use its directory
\path{Analysis/GeneralizedRiemannHypothesis}, at commit
\begin{center}
\small\texttt{7c24c5481fc0c8101b22b7e7cfab8321921c0faa}.
\end{center}
Its current article merges 194 earlier reports into 20 subject sections; its catalogue lists 1,200 theorem-like results. The README describes this as a merger of reports, without recomputing their mathematics. Our analysis used the current catalogue and source modules, followed by detailed examination of the parts relevant to these two extensions. Thirty-nine retrieved source files were checked against their pinned Git blob hashes. The accompanying provenance file records the inspected source set~\cite{ProveIt2026}.

The upstream result used by the article is RH and GRH for Dirichlet $L$-functions, at revision \texttt{72fead102f663d5c7a407964a13dc0cf9d977207} of \texttt{lsernau/grh-lean}. The exported statements use Mathlib's zeta and Dirichlet $L$-functions. Sutherland's verification page reports successful independent kernel and external-checker runs for this revision~\cite{Sutherland2026}. We did not replay that large verification. The first result below takes the source's Dirichlet-GRH theorem as input, and the second needs only finite linear algebra and probability.

The underlying research programme has two useful interfaces. On the analytic side, zero-location and explicit formulas convert oscillatory sums over zeros into prime-power coefficients. On the arithmetic and algebraic side, the source proof uses finite spanning estimates, tight frames, and Fourier estimates for matrix observations. The present extensions use these two interfaces separately. They do not transfer the zero-location theorem to another family of $L$-functions.

\subsection{First contribution: a quantitative finite-observation theorem}

The article's \texttt{li:thm:decoder} recovers $\Re\chi(p)$ from exact jumps of the Li--L\'evy process. Its \texttt{li:thm:grid} shows that every fixed finite grid has mutually absolutely continuous laws across characters. Question \texttt{li:q:paths} asks about discrete noisy observations. We answer its fixed-family, fixed-duration part.

For mesh $h$, increment-error bound $\eta$, and spectral bandwidth $T$, Theorem~\ref{thm:noisy-grid} constructs an estimator from positions alone with root-mean-square error bounded by
\begin{equation}\label{overview:basic}
C\left(\sqrt{\frac{\log T}{T}}
 +\frac{\log^2T}{T}+hT^3\log^2T+\eta T^2\log T\right)
\end{equation}
for fixed observation duration and sufficiently large $T$. The bound follows from the proof and the variance estimate in Theorem~\ref{thm:extrapolation}, with $m=1$. The error may be adversarial; there is no assumption of independent or mean-zero observation noise.

The conservative choice $T=\min(h^{-1/4},\eta^{-1/3})$ works for any $h\to0$ and $\eta\to0$, and also preserves the exact-jump Gaussian limiting covariance after multiplication by the window factor $\int w^2$. Thus one short interval contains enough arithmetic information for consistent finite-family identification even under vanishing measurement error. This is compatible with the impossibility of zero-error identification from one fixed finite grid.

Higher-order correction gives the achieved bound
\begin{equation}\label{overview:rates}
\operatorname{RMSE}\ \le C_m\max\left\{
h^{m/(4m+3)}\bigl(\log(1/h)\bigr)^{(3m+2)/(4m+3)},\quad
\eta^{1/5}\bigl(\log(1/\eta)\bigr)^{3/5}\right\}.
\end{equation}
Here $m$ is fixed, the second term is omitted for $\eta=0$, and $C_m$ also depends on the fixed arithmetic data. The mesh exponent approaches $1/4$ as the fixed order is increased. These are constructive upper bounds, not minimax claims. The conductor-uniform and optimal-information portions of the article's open question remain unresolved.

\subsection{Second contribution: the complete linear critical window}

An earlier delivered note, \emph{Recovering a Hidden Symplectic Form}, showed that $d=\binom n2-1$ is the threshold for recovering a hidden nondegenerate alternating form from random two-planes on which it vanishes. It proved success and failure across diverging additive windows, and a critical second-moment estimate. It did not determine the critical nullity distribution. This earlier note was inspected in its supplied working copy, with its hash recorded in the provenance file; it is more recent than the relevant merged article text.

The new Theorem~\ref{crit:main} gives the entire fixed-offset nullity law, with an explicit finite error bound. Its most direct consequence is
\[
\Pr\{N_{d+a}=\langle H\rangle\}
\longrightarrow \prod_{j=a+1}^{\infty}(1-q^{-j})
\qquad(t=d+a,\ a\ge0).
\]
Here $N_t$ is the space of forms satisfying the observed linear equations, so the event means that the hidden form is their only solution up to scale. For $a<0$ linear uniqueness is impossible. The result applies in every characteristic and uniformly in the hidden form. It also determines the optimal limiting test that retains only the solution-space dimension, and the distribution and all moments of the observation overhead.

The delicate point is universality. Each observation is a decomposable bivector and its coordinates are strongly dependent; in the planted experiment it lies in a prescribed symplectic hyperplane. The iid matrix law cannot simply be assumed. We prove that after an initial part of the sample has excluded forms of low rank, every subsequent decrease of the solution-space dimension has essentially the same probability as for a uniform random linear equation. The resulting coupling controls the whole final dimension process.

\subsection{Relation to prior work and the scope of novelty}

Small-time L\'evy moment expansions, high-frequency estimation, jump recovery, and Richardson extrapolation are established methods~\cite{Figueroa2008,Nickl2016,Gonzalez2021}. The arithmetic content here is the moving prime-frequency probe, its explicit sensitivity to noise, and its combination with a fixed observation horizon and a singular atomic zero measure. We give the required bounds for that setting rather than assuming that a fixed-test-function theorem covers it.

The rectangular iid nullity law is classical; Fulman and Goldstein give sharp total-variation results~\cite{FulmanGoldstein2015}. Low-rank counting and high-rank near-uniformity also appear in work on squares of random linear codes and products of codes~\cite{Cascudo2015,Randriambololona2015}. Cascudo and coauthors prove threshold and critical first-moment statements for quadratic evaluations. Their sampling model is different, and the inspected paper does not supply the fixed-offset planted symplectic nullity law proved here. The new finite-field content is the structured-row universality theorem and its explicit recovery, testing, and stopping-time consequences.

These are proved extensions beyond the inspected source statements and the preceding hidden-form note. The literature search establishes relevant antecedents and a scoped comparison; it does not certify worldwide publication priority. The report makes no claim that the new theorems have been independently peer reviewed or formally verified.

\begin{center}
\small
\begin{tabular}{@{}>{\raggedright\arraybackslash}p{.22\linewidth}>{\raggedright\arraybackslash}p{.34\linewidth}>{\raggedright\arraybackslash}p{.35\linewidth}@{}}
\toprule
Direction & Earlier result & Extension proved here\\
\midrule
Arithmetic paths & Exact jump decoder; overlapping fixed-grid laws & Noisy position estimator, fixed-horizon consistency, CLT and achieved mesh--precision rates\\[4pt]
Hidden symplectic form & Threshold across diverging sample gaps & Full fixed-offset nullity law; exact linear-recovery and dimension-test curves; stopping-time moments\\
\bottomrule
\end{tabular}
\end{center}
